% Part: first-order-logic
% Chapter: tableaux
% Section: quantifier-rules

\documentclass[../../../include/open-logic-section]{subfiles}

\begin{document}

\olfileid[id]{fol}{tab}{qrl}

\olsection{Aturan Kuantor}

\subsection{Aturan untuk $\lforall$}

\begin{defish}
\AxiomC{\sFmla{\True}{\lforall[x][!A(x)]}}
\RightLabel{\TRule{\True}{\lforall}}
\UnaryInfC{\sFmla{\True}{!A(t)}}
\DisplayProof
\hfill
\AxiomC{\sFmla{\False}{\lforall[x][!A(x)]}}
\RightLabel{\TRule{\False}{\lforall}}
\UnaryInfC{\sFmla{\False}{!A(a)}}
\DisplayProof
\end{defish}

Dalam \TRule{\True}{\lforall}, $t$ adalah suku tertutup (yaitu suku
tanpa variabel). Dalam \TRule{\False}{\lforall}, $a$ adalah
!!a{constant} yang tidak boleh muncul di mana pun pada cabang di atas
aturan \TRule{\False}{\lforall}. Kita menyebut $a$ sebagai
\emph{variabel eigen} dari inferensi
\TRule{\False}{\lforall}.
\footnote{Kita menggunakan istilah ``variabel eigen'' meskipun $a$
dalam aturan di atas adalah !!a{constant}. Alasannya bersifat historis.}

\subsection{Aturan untuk $\lexists$}

\begin{defish}
\AxiomC{\sFmla{\True}{\lexists[x][!A(x)]}}
\RightLabel{\TRule{\True}{\lexists}}
\UnaryInfC{\sFmla{\True}{!A(a)}}
\DisplayProof
\hfill
\AxiomC{\sFmla{\False}{\lexists[x][!A(x)]}}
\RightLabel{\TRule{\False}{\lexists}}
\UnaryInfC{\sFmla{\False}{!A(t)}}
\DisplayProof
\end{defish}

Sekali lagi, $t$ adalah suku tertutup, sedangkan $a$ adalah
!!a{constant} yang tidak muncul pada cabang di atas aturan
\TRule{\True}{\lexists}. Kita menyebut $a$ sebagai \emph{variabel
eigen} dari inferensi \TRule{\True}{\lexists}.

Syarat bahwa variabel eigen tidak muncul pada cabang di atas inferensi
\TRule{\False}{\lforall} atau \TRule{\True}{\lexists} disebut
\emph{syarat variabel eigen}.

\begin{explain}
Ingatlah konvensi bahwa ketika $!A$ adalah !!a{formula} dengan
!!{variable}~$x$ bebas, kita menandainya dengan menulis~$!A(x)$. Dalam
konteks yang sama, $!A(t)$ kemudian merupakan singkatan
untuk~$\Subst{!A}{t}{x}$. Jadi, kita juga dapat menuliskan aturan
\TRule{\False}{\lexists} sebagai:
\begin{prooftree}
  \AxiomC{\sFmla{\False}{\lexists[x][!A]}}
  \RightLabel{\TRule{\False}{\lexists}}
  \UnaryInfC{\sFmla{\False}{\Subst{!A}{t}{x}}}
\end{prooftree}
Perhatikan bahwa $t$ mungkin sudah muncul dalam~$!A$; misalnya, $!A$
mungkin berupa~$\Atom{\Obj P}{t,x}$. Dengan demikian, menyimpulkan
$\sFmla{\False}{\Atom{\Obj P}{t,t}}$ dari
$\sFmla{\False}{\lexists[x][\Atom{\Obj P}{t,x}]}$ merupakan penerapan
yang benar dari~\TRule{\False}{\lexists}. Akan tetapi, syarat variabel
eigen dalam \TRule{\False}{\lforall} dan~\TRule{\True}{\lexists}
mengharuskan agar !!{constant}~$a$ tidak muncul dalam~$!A$. Jadi, kita
tidak dapat menyimpulkan $\sFmla{\False}{\Atom{\Obj P}{a,a}}$ dengan
benar dari $\sFmla{\False}{\lforall[x][\Atom{\Obj P}{a,x}]}$ dengan
menggunakan~$\TRule{\False}{\lforall}$.
\end{explain}

\begin{explain}
Dalam \TRule{\True}{\lforall} dan \TRule{\False}{\lexists}, suku~$t$
boleh berupa sebarang suku tertutup. Sebaliknya, dalam aturan
\TRule{\True}{\lexists} dan \TRule{\False}{\lforall}, syarat variabel
eigen mengharuskan agar !!{constant}~$a$ tidak muncul di mana pun pada
cabang di atas inferensi yang bersangkutan. Syarat ini diperlukan untuk
memastikan bahwa sistem tersebut sahih. Tanpa syarat ini, berikut akan
menjadi !!{tableau} tertutup untuk
$\lexists[x][\formula{A}(x)] \lif \lforall[x][\formula{A}(x)]$:
\begin{center}
\begin{tableau}{}
  [\sFmla{\False}{\lexists[x][\formula{A}(x)] \lif \lforall[x][\formula{A}(x)]}, just=\TAss
    [\sFmla{\True}{\lexists[x][\formula{A}(x)]},
      just={\TRule{\False}{\lif}[1]}
      [\sFmla{\False}{\lforall[x][\formula{A}(x)]},
        just={\TRule{\False}{\lif}[1]}
        [\sFmla{\True}{\formula{A}(a)},
          just={\TRule{\True}{\lexists}[2]}
          [\sFmla{\False}{\formula{A}(a)},
            just={\TRule{\False}{\lforall}[3]}, close]
        ]
      ]
    ]
  ]
\end{tableau}
\end{center}
Namun, $\lexists[x][\formula{A}(x)] \lif
\lforall[x][\formula{A}(x)]$ tidak valid.
\end{explain}

\end{document}
